% TOP_PROBLEM: {"id": 2305003, "title": "Research Problems in Function Theory — Problem 5.3", "category_id": 8, "category": {"id": 8, "name": "analysis", "display_name": "Analysis", "description": "Limits, continuity, calculus, and function theory.", "slug": "analysis", "order_index": 8, "created_at": "2026-07-31T15:26:25.671Z"}, "status": "open"}
% FIRST_POSTED: null
% ENTRY_KIND: statement_only
% Imported from: batch13.tex; UnsolvedMath ID 2305003
\documentclass[11pt]{book}
\usepackage{fontspec}
\setmainfont{Latin Modern Roman}
\setsansfont{Latin Modern Sans}
\usepackage{amsmath,amssymb,mathtools}
\usepackage{unicode-math}
\setmathfont{Latin Modern Math}
\usepackage{xurl}
\usepackage[hidelinks]{hyperref}
\newcommand{\sref}[2]{\hyperlink{source:#1:#2}{[S#2]}}
\newcommand{\cataloguescope}[1]{\paragraph{Mathematical question.}#1}
\begin{document}
% UnsolvedMath ID 2305003; source catalogue code AMR-022-5003
\setcounter{section}{2305002}
\section{Research Problems in Function Theory — Problem 5.3}\label{2305003}
{\small\sffamily UnsolvedMath ID 2305003 · Analysis}
\subsection{Definitions and mathematical statement}

\paragraph{Shared notation.}$\mathbb N=\{1,2,\ldots\}$, and $\mathbb Z,\mathbb Q,\mathbb R,\mathbb C$ denote the integers, rationals, reals and complex numbers. Any other conventions are specified locally.


Let $\mathbb D=\{z\in\mathbb C:|z|<1\}$ and let $f(z)=\sum_{n=0}^\infty a_nz^n$ be holomorphic in $\mathbb D$. For $\lambda>0$ and $0<r<1$ define the integral mean
\[I_\lambda(r,f)=\left(\frac1{2\pi}\int_0^{2\pi}|f(re^{i\theta})|^\lambda\,d\theta\right)^{1/\lambda}.
\]
The coefficient estimates refer to $n\to\infty$, and mean estimates to $r\uparrow1$. Constants in $O(\cdot)$ may depend on $f$ and the indicated exponent, but not on $n$ or $r$. A function is \emph{weakly univalent} here if, for every radius $R>0$, one of the following alternatives holds: every value $w$ on $|w|=R$ is assumed exactly once in $\mathbb D$, or at least one value $w_R$ on that circle is omitted throughout $\mathbb D$. Preimages in the first alternative are counted with their holomorphic multiplicities. This condition is imposed directly, without requiring an independently chosen omitted-value sequence.
The earlier estimates recorded in the problem are $I_\lambda(r,f)=O((1-r)^{-2})$ for $\lambda>0$ and $|a_n|=O(n^2)$.
\paragraph{Statement recorded in the supplied source.}
For weakly univalent $f$, obtain sharper general bounds for the integral means $I_\lambda(r,f)$ and for $|a_n|$. In particular, decide whether the stronger bounds
\[I_1(r,f)=O\bigl((1-r)^{-1}\bigr),\qquad |a_n|=O(n)
\]
hold, and seek improvements of the quadratic estimates even independently of these full-strength targets. \sref{2305003}{2}
\subsection{Short English statement}
Sharpen the integral-mean and coefficient bounds for weakly univalent functions, with the first-mean boundary-distance bound and linear coefficients as the stronger targets.
\subsection{Sources}
\begin{description}
\item[\hypertarget{source:2305003:1}{S1}] ULAM.AI and the UnsolvedMath Contributors, \emph{UnsolvedMath: A Curated Collection of Open Mathematics Problems}, record 2305003 (AMR-022-5003). CC BY 4.0. \url{https://huggingface.co/datasets/ulamai/UnsolvedMath}.
\item[\hypertarget{source:2305003:2}{S2}] Walter K. Hayman and Edward F. Lingham, \emph{Research Problems in Function Theory}, arXiv:1809.07200, 2018. Chapter 5 introduction, equations (5.1)--(5.6), Problem 5.3, and complete Updates 5.1--5.3. \url{https://arxiv.org/pdf/1809.07200}.
\end{description}
\label{end:2305003}
\end{document}
